\documentclass[a5paper,10pt]{article}

% https://github.com/InDevRus/filippov-solutions

% Нумерация страниц
\pagenumbering{gobble}

% Настройка полей
\usepackage{geometry}
\geometry{left=1cm}
\geometry{right=1cm}
\geometry{top=1cm}
\geometry{bottom=1.5cm}

% Вставка таблицы
\usepackage{array}
\usepackage{tabularx}

% Инструменты выравнивания формул
\usepackage{amsmath}
\usepackage{mathtools}

% Диакритика в математике
\usepackage{amsfonts}

% Вычисления размеров на ходу
\usepackage{calc}

% Удобные записи производных
\usepackage[italic=true]{derivative}

% Настройки сносок
\usepackage{enotez}

% Длинные знаки векторов
\usepackage{anyfontsize}
\usepackage[b]{esvect}

% Вставка картинок
\usepackage{graphicx}

% Вставка красной строки
\usepackage{indentfirst}
\setlength{\parindent}{1em}

% Условия в командах
\usepackage{xifthen}

% Луа-скрипты
\usepackage{luacode}

% Настройка команд отдельно в математическом и текстовом режимах
\usepackage{mathcommand}

% Для повторения LaTeX кода
\usepackage{multido}

% Конфигурация отступов формул
\delimitershortfall-1sp
\usepackage{mleftright}
\mleftright

% Растягивание объектов
\usepackage{relsize}

% Обтекание текстом
\usepackage{wrapfig}
\usepackage{subcaption}
\usepackage{floatflt}

% Поддержка ввода кириллицы
\usepackage[english, russian]{babel}

% Настройка корневой позиции для компилятора
\usepackage{import}
\providecommand*{\ProjectRootPath}{..}

\import{\ProjectRootPath/preambles/macros/}{theorems}

\import{\ProjectRootPath/preambles/fonts}{font_setup}

% Знак градуса
\usepackage{siunitx}

\import{\ProjectRootPath/preambles/macros/}{operators}
\import{\ProjectRootPath/preambles/macros/}{redefinitions}

\import{\ProjectRootPath/preambles/fonts}{letters_setup}
\import{\ProjectRootPath/preambles/fonts/adjustments}{custom_superscripts}

\import{\ProjectRootPath/preambles/custom}{formatting}
\import{\ProjectRootPath/preambles/custom}{frames}
\import{\ProjectRootPath/preambles/custom}{list_setup}

% TODO: перевести команды на современный стиль объявления
\input{../preambles/plotting}

\begin{document}

$y\`\br{x} = \dfrac {2x - y} {x - y}$. Это уравнение задаёт траектории, например,
системы $\tsys{\dot{x} = x - y \\ \dot{y} = 2x - y}$.
$A = \begin{pmatrix} 1 & -1 \\ 2 & -1 \end{pmatrix}$.
$\det\br{\lambda E - A} = \br{\lambda - 1}\br{\lambda + 1} + 2 = \pow{\lambda}{2} + 1$.
$\sub{\lambda}{1} = -i$, $\sub{\lambda}{2} = i$.

Итак, особая точка $\br{0; 0}$ -- это <<центр>>.

Это означает, что решение дифференциального уравнения можно получить даже в явном виде.
$y\` = \dfrac {2 - z} {1 - z}$, где
$y = z\br{x} \cdot x$. \vspace{1mm}
\begin{align*}
    \dfrac {z - 1} {\pow{z}{2} - 2z + 2} z\`\br{x} + \dfrac {1} {x} = 0 &&
    \dfrac {1} {2} \ln\vbr{\pow{z}{2} - 2z + 2} + \ln\vbr{x} = C
\end{align*}
$\br{\br{z - 1}^2 + 1}\, \pow{x}{2} = \pow{C}{2}$.
$ - x}^2 + \pow{x}{2} = \pow{C}{2}$.
$\pow{y}{2} + 2\pow{x}{2} - 2xy = \pow{C}{2}$.
Эта квадратичная форма задаёт эллипсы, в чём несложно убедиться, приведя эту форму к
каноническому виду, для чего достаточно сделать поворот на угол
$\alpha = \arctg\br{\dfrac {1 + \sqrt{5}} {2}}$.
Подстановка 
$\tsys{
    x = \cos\br{\alpha} u - \sin\br{\alpha} v 
    \\ y = \sin\br{\alpha} u + \cos\br{\alpha} v
}$
приводит к форме 
$\dfrac {\pow{u}{2}} {\pow{a}{2}} + \dfrac {\pow{v}{2}} {\pow{b}{2}} 
= \pow{C}{2}$, 
где $a = \dfrac {1 + \sqrt{5}} {2}$, $b = \dfrac {1 - \sqrt{5}} {2}$.

\begin{floatingfigure}{0.5\textwidth}
    \begin{tikzpicture} 
        \begin{axis}[
                width = \textwidth,
                height = 9.5cm,
                ymin = -3.4,
                ymax = 3.4,
                xmin = -2.4,
                xmax = 2.4, 
                xtick = {-3, -2, ..., 3},
                ytick = {-5, -4, ..., 5},
                minor tick num = 1,
            ],
            \newcommand{\ALPHAA}{(1 + sqrt(5)) / 2}
            \newcommand{\ALPHAB}{(1 - sqrt(5)) / 2}
            
            \addplot[domain = -5 : 5, dashed]
            {\ALPHAA * x};
            
            \draw [thick]  (0, 0) ellipse [
                rotate = atan(\ALPHAA),
                x radius = \ALPHAA * 0.25,
                y radius = \ALPHAB * 0.25,
            ];
            \draw [thick]  (0, 0) ellipse [
                rotate = atan(\ALPHAA),
                x radius = \ALPHAA * 0.5,
                y radius = \ALPHAB * 0.5,
            ];
            \draw [thick]  (0, 0) ellipse [
                rotate = atan(\ALPHAA),
                x radius = \ALPHAA,
                y radius = \ALPHAB,
            ];
            \draw [thick]  (0, 0) ellipse [
                rotate = atan(\ALPHAA),
                x radius = \ALPHAA * 1.5,
                y radius = \ALPHAB * 1.5,
            ];
            \draw [thick]  (0, 0) ellipse [
                rotate = atan(\ALPHAA),
                x radius = \ALPHAA * 2,
                y radius = \ALPHAB * 2,
            ];
        \end{axis}
    \end{tikzpicture}
\end{floatingfigure}

Соответственно, ось эллипса есть \linebreak прямая
$y = \tg\br{\alpha} x $, то есть $y = \dfrac {1 + \sqrt{5}} {2}\, x$.
Эта прямая обозначена пунктирной \linebreak линией.

\end{document}
